% GENERATED FILE. Edit canonical lessons, then run npm run generate:books.
% canonical-source-hash: fnv1a64:0061d91f334a4b61
% canonical-lessons: BN-C45-khushi, BN-C45-daktar, BN-C45-baba, BN-C45-chhele, BN-C45-meye

\chapter{Happy, Doctor, Father, Son, Daughter}
\label{ch:bn-happy-doctor-father-son-daughter}
\hlchaptermodality{45}

\begin{chapteropening}
\textbf{By the end of this chapter:} \emph{I can say happy, a doctor, a father, a son and a daughter.}

Complete the last lesson of chapter 45: I can say happy, a doctor, a father, a son and a daughter.
\end{chapteropening}

\section[\texorpdfstring{\bn{খুশি} (khushi)}{খুশি (khushi)}]{\bn{খুশি} (khushi) — happy}

\label{lesson:BN-C45-khushi}



\begin{quote}
Before the new one: say the Bengali for pain, then the Bengali for hunger.
\end{quote}

\subsection*{You'll want to know: \bn{খুশি}}
\textbf{\bn{খুশি}} — \emph{khushi} — \textquotedblleft{}happy\textquotedblright{}.

\begin{cousinweb}
\textbf{\bn{খুশি}} (\emph{khushi}) is a loan, usually traced to Persian \emph{khushī}, \textquotedblleft{}happiness, joy\textquotedblright{}. It is built on Persian \emph{khush}, \textquotedblleft{}happy\textquotedblright{}.
\end{cousinweb}

\begin{grammarlens}[title={What you've built}]
One more word for this chapter.
\end{grammarlens}

\subsection*{Your turn}
\begin{itemize}
\raggedright
  \item \emph{Say it:} \emph{khushi}
  \item \emph{Say it:} \emph{khushi}, once more
  \item \emph{Say it:} say \emph{khushi}
  \item \emph{Recall:} say the Bengali for pain, then the Bengali for hunger, then say \emph{khushi} again
\end{itemize}

\begin{tcolorbox}[breakable,colback=teal!4,colframe=teal!35!black,title={Before you move on}]
What does \bn{খুশি} mean? (\textquotedblleft{}Happy\textquotedblright{}.) Say it once more.
\end{tcolorbox}
\section[\texorpdfstring{\bn{ডাক্তার} (ḍāktār)}{ডাক্তার (ḍāktār)}]{\bn{ডাক্তার} (ḍāktār) — a doctor}

\label{lesson:BN-C45-daktar}



\begin{quote}
Before the new one: say the Bengali for hunger, then the Bengali for happy.
\end{quote}

\subsection*{You'll want to know: \bn{ডাক্তার}}
\textbf{\bn{ডাক্তার}} — \emph{ḍāktār} — \textquotedblleft{}a doctor\textquotedblright{}.

\begin{cousinweb}
\textbf{\bn{ডাক্তার}} (\emph{ḍāktār}) is a loan from English \emph{doctor}.
\end{cousinweb}

\begin{grammarlens}[title={What you've built}]
One more word for this chapter.
\end{grammarlens}

\subsection*{Your turn}
\begin{itemize}
\raggedright
  \item \emph{Say it:} \emph{ḍāktār}
  \item \emph{Say it:} \emph{ḍāktār}, once more
  \item \emph{Say it:} say \emph{ḍāktār}
  \item \emph{Recall:} say the Bengali for hunger, then the Bengali for happy, then say \emph{ḍāktār} again
\end{itemize}

\begin{tcolorbox}[breakable,colback=teal!4,colframe=teal!35!black,title={Before you move on}]
What does \bn{ডাক্তার} mean? (\textquotedblleft{}A doctor\textquotedblright{}.) Say it once more.
\end{tcolorbox}
\section[\texorpdfstring{\bn{বাবা} (bābā)}{বাবা (bābā)}]{\bn{বাবা} (bābā) — a father}

\label{lesson:BN-C45-baba}



\begin{quote}
Before the new one: say the Bengali for happy, then the Bengali for a doctor.
\end{quote}

\subsection*{You'll want to know: \bn{বাবা}}
\textbf{\bn{বাবা}} — \emph{bābā} — \textquotedblleft{}a father\textquotedblright{}.

\begin{cousinweb}
\textbf{\bn{বাবা}} (\emph{bābā}) is a tadbhava, a word inherited and reshaped by sound change. It comes from \emph{*bābba}, \textquotedblleft{}father\textquotedblright{}, a form CDIAL reconstructs from its descendants (CDIAL 9209). CDIAL calls these nursery words.
\end{cousinweb}

\begin{grammarlens}[title={What you've built}]
One more word for this chapter.
\end{grammarlens}

\subsection*{Your turn}
\begin{itemize}
\raggedright
  \item \emph{Say it:} \emph{bābā}
  \item \emph{Say it:} \emph{bābā}, once more
  \item \emph{Say it:} say \emph{bābā}
  \item \emph{Recall:} say the Bengali for happy, then the Bengali for a doctor, then say \emph{bābā} again
\end{itemize}

\begin{tcolorbox}[breakable,colback=teal!4,colframe=teal!35!black,title={Before you move on}]
What does \bn{বাবা} mean? (\textquotedblleft{}A father\textquotedblright{}.) Say it once more.
\end{tcolorbox}
\section[\texorpdfstring{\bn{ছেলে} (chhele)}{ছেলে (chhele)}]{\bn{ছেলে} (chhele) — a son; a boy}

\label{lesson:BN-C45-chhele}



\begin{quote}
Before the new one: say the Bengali for a doctor, then the Bengali for a father.
\end{quote}

\subsection*{You'll want to know: \bn{ছেলে}}
\textbf{\bn{ছেলে}} — \emph{chhele} — \textquotedblleft{}a son; a boy\textquotedblright{}.

\begin{cousinweb}
CDIAL, in its addenda, lists \textbf{\bn{ছেলে}} (\emph{chhele}), \textquotedblleft{}boy\textquotedblright{}, under Sanskrit \emph{chagala}, \textquotedblleft{}goat\textquotedblright{}, beside the form \emph{chēlaka}, \textquotedblleft{}goat\textquotedblright{} (CDIAL 4963): the word for a boy seems to have begun as a word for a young goat.
\end{cousinweb}

\begin{grammarlens}[title={What you've built}]
One more word for this chapter.
\end{grammarlens}

\subsection*{Your turn}
\begin{itemize}
\raggedright
  \item \emph{Say it:} \emph{chhele}
  \item \emph{Say it:} \emph{chhele}, once more
  \item \emph{Say it:} say \emph{chhele}
  \item \emph{Recall:} say the Bengali for a doctor, then the Bengali for a father, then say \emph{chhele} again
\end{itemize}

\begin{tcolorbox}[breakable,colback=teal!4,colframe=teal!35!black,title={Before you move on}]
What does \bn{ছেলে} mean? (\textquotedblleft{}A son; a boy\textquotedblright{}.) Say it once more.
\end{tcolorbox}
\section[\texorpdfstring{\bn{মেয়ে} (meye)}{মেয়ে (meye)}]{\bn{মেয়ে} (meye) — a daughter; a girl}

\label{lesson:BN-C45-meye}



\begin{quote}
Before the new one: say the Bengali for a father, then the Bengali for a son; a boy.
\end{quote}

\subsection*{You'll want to know: \bn{মেয়ে}}
\textbf{\bn{মেয়ে}} — \emph{meye} — \textquotedblleft{}a daughter; a girl\textquotedblright{}.

\begin{cousinweb}
CDIAL lists \textbf{\bn{মেয়ে}} (\emph{meye}), \textquotedblleft{}girl\textquotedblright{}, under Sanskrit \emph{mātṛ}, \textquotedblleft{}mother\textquotedblright{}, the same entry that gives Bengali \emph{mā}, \textquotedblleft{}mother\textquotedblright{} (CDIAL 10016).
\end{cousinweb}

\begin{grammarlens}[title={What you've built}]
One more word for this chapter.
\end{grammarlens}

\subsection*{Your turn}
\begin{itemize}
\raggedright
  \item \emph{Say it:} \emph{meye}
  \item \emph{Say it:} \emph{meye}, once more
  \item \emph{Say it:} say \emph{meye}
  \item \emph{Recall:} say the Bengali for a father, then the Bengali for a son; a boy, then say \emph{meye} again
\end{itemize}

\begin{tcolorbox}[breakable,colback=teal!4,colframe=teal!35!black,title={Before you move on}]
What does \bn{মেয়ে} mean? (\textquotedblleft{}A daughter; a girl\textquotedblright{}.) Say it once more.
\end{tcolorbox}
